\chapter{Đọc thêm 1: Giáo dục kỹ thuật gắn với tư duy phản biện và giải quyết vấn đề (Ismail)}
\label{ch:M5L4DocThem1}

\section{Thông tin nguồn}

\begin{itemize}
  \item \textbf{Nguồn:} Ismail, M. \emph{IOP Conf. Ser.: Mater. Sci. Eng.}, 697, 012017, 2019.
  \item \textbf{Phần cần đọc:} Toàn bộ bài báo
  \item \textbf{Ghi chú:} bản chuyển ngữ tiếng Việt súc tích, bám cấu trúc nguồn (giữ nguyên tiêu đề, công thức, số liệu, bảng, chú thích và danh mục tài liệu tham khảo).
\end{itemize}

IOP Conference Series: Materials Science and Engineering, 697 (2019) 012017, ICE\&ICIE 2019. Bài báo truy cập mở (PAPER • OPEN ACCESS). doi:10.1088/1757-899X/697/1/012017

W Omar Ali Saifuddin Wan Ismail (1*), Noraini Hamzah (2), Ireana Yusra Abdul Fatah (1) và Azry Khoiry Muhammad (2)

\begin{enumerate}
\def\labelenumi{\arabic{enumi}.}
\item
  Faculty of Innovative Design and Technology, Universiti Sultan Zainal Abidin, Gong Badak Campus, 21300 Kuala Terengganu, Malaysia
\item
  Faculty of Engineering and Built Environment, Universiti Kebangsaan Malaysia, 43600 Bangi, Malaysia
\end{enumerate}

Email: woasaifuddin@unisza.edu.my

\textbf{Tóm tắt.} Giáo dục kỹ thuật là học phần nhấn mạnh các kỹ năng mềm liên quan đến tư duy phản biện (critical thinking) và khả năng giải quyết vấn đề. Kỹ năng tư duy sáng tạo và phản biện (CCTS) đã được đưa vào và đặc biệt chú trọng trong hệ thống giáo dục từ bậc tiểu học. Gần đây, toàn bộ đề thi ở tiểu học, trung học và đại học công lập đều định hướng theo nguyên tắc Kỹ năng tư duy bậc cao (HOTS) trong khuôn khổ RMK 11 (Chuyển đổi hệ thống giáo dục). Mục tiêu của chương trình là đào tạo sinh viên có năng lực giải quyết vấn đề tốt hơn nhờ vận dụng hiệu quả nhiều dạng kiến thức đã học về lý thuyết. Nghiên cứu này xem xét mức độ cần thiết của các yếu tố tư duy phản biện và kỹ năng giải quyết vấn đề để đào tạo ra những kỹ sư có sức cạnh tranh, từ góc nhìn của ngành công nghiệp. Phản hồi được thu thập bằng phương pháp định lượng thông qua bộ bảng hỏi, với tổng cộng 300 người tham gia. Nghiên cứu được kỳ vọng sẽ là tài liệu định hướng cho các kỹ sư cơ khí tương lai trước những yêu cầu của ngành công nghiệp chủ đạo.

\section{1. Giới thiệu}

Ở thế kỷ 19, tư duy phản biện được hiểu là khả năng phân tích thông tin, xác định mức độ liên quan của dữ liệu thu thập và diễn giải để giải quyết vấn đề {[}1-5{]}. Đến thế kỷ 20, tư duy phản biện được xem là sự thể hiện suy nghĩ của con người qua việc tư duy về các hoạt động mình thực hiện {[}6{]}, ở cấp độ cao hơn (kỹ năng tư duy bậc cao), bao gồm quá trình phân tích, đánh giá, công bằng và phản tư {[}7,8{]}. Tư duy phản biện cũng hỗ trợ sự phát triển nhận thức của sinh viên khi họ cố gắng thu thập càng nhiều thông tin càng tốt để giải quyết vấn đề, tổ chức thông tin và nhìn nhận vấn đề từ nhiều góc độ khác nhau {[}9,10{]}. Tư duy phản biện là kiểu tư duy có chủ đích, hướng đến mục tiêu, trong quá trình ra quyết định dựa trên bằng chứng của việc giải quyết vấn đề khoa học, chứ không phải chỉ phỏng đoán {[}11,12{]}. Từ đó, tư duy phản biện và phương pháp phán đoán là yêu cầu đối với kỹ sư ở cấp độ chuyên nghiệp để giải quyết phần lớn các vấn đề trong công việc {[}13{]}, và là một trong những điều kiện để có được việc làm tốt, đặc biệt trong lĩnh vực kỹ thuật {[}14-16{]}.

Quá trình giải quyết vấn đề là yêu cầu thiết yếu đối với mọi kỹ sư {[}17{]} và cũng quan trọng với mọi vị trí chuyên môn khác {[}18{]}. Do đó, cần nuôi dưỡng xu hướng tư duy phản biện, xem đó là thành phần thiết yếu trong sự phát triển quá trình tư duy và cách tiếp cận giải quyết vấn đề, từ đó hình thành chuyên môn {[}19{]}. Việc tập trung vào kỹ năng giải quyết vấn đề đã được nêu ở mọi cấp học, nhất là ở các cơ sở giáo dục đại học, nhưng kết quả nghiên cứu cho thấy khả năng giải quyết vấn đề của sinh viên kỹ thuật vẫn ở mức thấp, bất chấp nhiều nỗ lực của giảng viên trong chương trình cử nhân bốn năm {[}20{]}. Sinh viên kỹ thuật là những người cần tư duy phản biện nhiều hơn sinh viên các ngành khoa học xã hội, khoa học cơ bản và nhân văn {[}21-23{]}.

Nhu cầu về tư duy phản biện và kỹ năng giải quyết vấn đề sẽ tiếp tục tăng trong những năm tới, do nhu cầu cấp thiết về những hiểu biết mới về cách con người xử lý sự phức tạp và bất định {[}24{]}. Chương trình i-think bắt đầu được triển khai năm 2014 như một công cụ tư duy nhằm thúc đẩy tư duy bậc cao, nâng cao và bồi dưỡng kỹ năng tư duy cho sinh viên {[}25{]}. Mọi nhà tuyển dụng và tổ chức trong ngành đều đang cố gắng tuyển dụng và đào tạo kỹ sư thuộc mọi chuyên ngành dựa trên kỹ năng giải quyết vấn đề và tư duy phản biện, kết hợp với kỹ năng giao tiếp {[}26, 27{]}.

\section{2. Thực nghiệm (Experimental)}

Phương pháp phân tích phân biệt (Discriminant Analysis) được dùng để xử lý dữ liệu. Có tổng cộng 300 người tham gia trả lời khảo sát, đều đến từ doanh nghiệp, gồm nhóm quản lý và hành chính, kỹ sư cao cấp, kỹ thuật viên cao cấp và kỹ sư. Các mục hỏi phản ánh mức độ đáp ứng yêu cầu đối với kỹ năng tư duy phản biện và giải quyết vấn đề theo nhu cầu hiện nay của doanh nghiệp. Từ tổng điểm các mục này, nghiên cứu tính tỷ lệ phần trăm và giá trị trung bình cho kỹ năng tư duy phản biện và giải quyết vấn đề, được trình bày trong các bảng.

\subsection{2.1 Mô hình kỹ năng mềm của cơ sở giáo dục đại học}

Theo hình 1, Mô hình các yếu tố kỹ năng mềm do Bộ Giáo dục (MoE) Malaysia đưa ra, là nguồn chính tạo nên thành công và kết quả học tập của sinh viên ở một cấp độ nhất định.

Các yếu tố trong mô hình: Giao tiếp; Làm việc nhóm; Lãnh đạo; Học tập liên tục và quản lý thông tin; Tư duy phản biện và kỹ năng giải quyết vấn đề; Tinh thần khởi nghiệp; Đạo đức và chuẩn mực nghề nghiệp.

Hình 1. Mô hình kỹ năng mềm (MoHE).

Trong nghiên cứu này, các tác giả sử dụng mô hình gồm bảy yếu tố nêu trên, trong đó một yếu tố quan trọng là tư duy phản biện và kỹ năng giải quyết vấn đề. Phần sau sẽ trình bày tầm quan trọng của việc rèn luyện kỹ năng tư duy cho sinh viên ở cả đại học công lập và tư thục, bên cạnh giáo dục trước đó. Vì cải cách giáo dục đang tập trung vào quá trình tư duy phản biện, gồm cả việc giải quyết vấn đề {[}28-30{]} như điểm khởi đầu về nhận thức {[}31{]}, quá trình phát triển tiếp tục tạo tác động tích cực đối với môi trường việc làm.

Mô hình Hình thành Tư duy Hạng nhất (Model of Formation of First Class Mentality) do Bộ Giáo dục (MOE) Malaysia xây dựng, như hình 2.

Ba cấp độ của mô hình (hình 2):

\begin{itemize}
\item
  Cấp 1: Tiềm năng thay đổi; Bắt đầu suy nghĩ nghiêm túc; Có thể chấp nhận thách thức; Bắt đầu đánh giá năng lực bản thân.
\item
  Cấp 2: Tư duy phản biện; Chấp nhận thay đổi; Cởi mở; Khả năng tư duy.
\item
  Cấp 3: Think tank (tập hợp trí tuệ); Nhà lãnh đạo có tầm nhìn; Có trí tuệ và tận tâm; Niềm tin; Dám chấp nhận rủi ro.
\end{itemize}

Hình 2. Mô hình hình thành tư duy hạng nhất.

Tầm quan trọng và nhu cầu đối với kỹ năng tư duy được thể hiện rõ trong mô hình năng lực tư duy. Yếu tố Tư duy (Think) được xem là phần nền tảng của mọi giai đoạn trong quá trình gọi là tư duy sáng tạo và phản biện, nhằm đạt mục tiêu cuối cùng của giáo dục kỹ thuật ở đại học. Các tiêu chí của mô hình được chi tiết hóa khi xây dựng bảng hỏi. Mỗi cấp độ nhấn mạnh tầm quan trọng của kỹ năng tư duy theo những cách khác nhau trong giải quyết vấn đề và ra quyết định.

Hiện nay, khả năng nhận diện và giải quyết vấn đề ở sinh viên kỹ thuật tại Malaysia được đòi hỏi rất cao. Yếu tố thiết yếu của thực hành tư duy phản biện là phản tư (reflection) {[}32{]}, giúp cá nhân tư duy sắc bén, phân tích khi giải quyết mọi vấn đề, đồng thời hình thành thái độ tích cực khi đối mặt với các vấn đề {[}33{]}. Khi áp dụng tư duy phản biện, sinh viên cần khai thác tính sáng tạo bằng cách xác định vấn đề bằng một công cụ sáng tạo phổ biến (5W+1H: cái gì, ai, ở đâu, khi nào, tại sao và như thế nào) để nhận diện và giải quyết vấn đề {[}34{]} với nhiều phương án đa dạng.

Trong bối cảnh hiện nay, sinh viên tốt nghiệp gặp khó khăn khi tìm việc do thiếu khả năng tư duy phản biện, ngoài kiến thức và kỹ năng mềm {[}35{]}. Một biện pháp tốt là xây dựng mô hình giảng dạy và đánh giá theo tiêu chí Kỹ năng tư duy bậc cao (HOTS), công khai khuyến khích sinh viên dùng tư duy sáng tạo và phản biện khi giải quyết vấn đề {[}36, 37{]}. Sinh viên kỹ thuật không chỉ cần trình độ và kỹ năng chuyên môn để đảm nhiệm công việc hằng ngày; nếu thiếu kỹ năng mềm, nhà tuyển dụng ít có khả năng nhận họ vào vị trí chuyên nghiệp, quản lý hay thậm chí đối tác kinh doanh {[}38{]}. Với sinh viên kỹ thuật hiện đại, kỹ năng mềm hay kỹ năng chung là yếu tố phải được thành thạo và lồng ghép vào mọi khía cạnh kỹ năng nếu muốn hoàn thành công việc thành công, đồng thời góp phần nâng cao năng lực và sự phát triển của công ty {[}39, 40{]}.

\section{3. Kết quả và thảo luận}

\subsection{3.1 Phân tích thống kê mô tả nhân khẩu học}

Bảng 1 cho thấy biểu đồ tần suất về vị trí hiện tại của người trả lời và thống kê mô tả nhân khẩu học: họ có hơn năm năm kinh nghiệm trong ngành, đặc biệt trong lĩnh vực kỹ thuật cơ khí và các phòng ban kỹ thuật. Năm năm kinh nghiệm được xem là mốc khởi đầu của giai đoạn trưởng thành trong con đường nghề nghiệp.

Bảng 1. Thống kê mô tả nhân khẩu học.

{\small
\begin{longtable}[]{@{}llll@{}}
\toprule\noalign{}
STT & Vị trí hiện tại & Tần suất & Tỷ lệ phần trăm \\
\midrule\noalign{}
\endhead
\bottomrule\noalign{}
\endlastfoot
1 & Quản lý và hành chính & 100 & 33.33 \\
2 & Kỹ sư cao cấp & 72 & 24.00 \\
3 & Kỹ thuật viên cao cấp & 38 & 12.67 \\
4 & Kỹ sư & 90 & 30.00 \\
\end{longtable}
}

Bảng 2 cho thấy biểu đồ tần suất về giới tính và thống kê mô tả nhân khẩu học: trong số người trả lời có 240 nam và 60 nữ. Mẫu gồm nhân viên các ngành công nghiệp nặng (HICOM), vừa và nhỏ (SME) trên khắp bán đảo Malaysia. Nhìn chung, nhân sự ngành kỹ thuật cơ khí chủ yếu là nam giới.

Bảng 2. Thống kê mô tả nhân khẩu học.

{\small
\begin{longtable}[]{@{}llll@{}}
\toprule\noalign{}
STT & Giới tính & Tần suất & Tỷ lệ phần trăm \\
\midrule\noalign{}
\endhead
\bottomrule\noalign{}
\endlastfoot
1 & Nam & 240 & 80.00 \\
2 & Nữ & 60 & 20.00 \\
\end{longtable}
}

\subsection{3.2 Kết quả của bảy mục hỏi trong bảng hỏi}

Bảng hỏi được phát cho người trả lời để thu thập câu trả lời có tính xây dựng. Các câu hỏi được chọn trong bảng hỏi như bảng 3.

Bảng 3. Các mục của bảng hỏi.

{\small
\begin{longtable}[]{@{}
  >{\raggedright\arraybackslash}p{(\columnwidth - 2\tabcolsep) * \real{0.5000}}
  >{\raggedright\arraybackslash}p{(\columnwidth - 2\tabcolsep) * \real{0.5000}}@{}}
\toprule\noalign{}
\begin{minipage}[b]{\linewidth}\raggedright
Câu
\end{minipage} & \begin{minipage}[b]{\linewidth}\raggedright
Nội dung câu hỏi
\end{minipage} \\
\midrule\noalign{}
\endhead
\bottomrule\noalign{}
\endlastfoot
Q1 & Khả năng nhận diện và phân tích vấn đề trong tình huống phức tạp, mơ hồ và đưa ra đánh giá có căn cứ. \\
\end{longtable}
}

{\small
\begin{longtable}[]{@{}
  >{\raggedright\arraybackslash}p{(\columnwidth - 2\tabcolsep) * \real{0.1818}}
  >{\raggedright\arraybackslash}p{(\columnwidth - 2\tabcolsep) * \real{0.8182}}@{}}
\toprule\noalign{}
\begin{minipage}[b]{\linewidth}\raggedright
Mã
\end{minipage} & \begin{minipage}[b]{\linewidth}\raggedright
Nội dung câu hỏi
\end{minipage} \\
\midrule\noalign{}
\endhead
\bottomrule\noalign{}
\endlastfoot
Q2 & Khả năng mở rộng và cải thiện các kỹ năng tư duy như giải thích, phân tích và đánh giá trong thảo luận. \\
Q3 & Khả năng tìm kiếm ý tưởng và tìm các giải pháp thay thế. \\
Q4 & Khả năng tư duy sáng tạo, vượt ra ngoài lối mòn. \\
Q5 & Khả năng ra quyết định dựa trên bằng chứng vững chắc. \\
Q6 & Khả năng kiên trì và hoàn thành đầy đủ trách nhiệm được giao. \\
Q7 & Khả năng hiểu và thích nghi với văn hóa cộng đồng và môi trường làm việc mới. \\
\end{longtable}
}

Bảng 4 trình bày kết quả của bảng hỏi bảy mục trong nghiên cứu này. Theo kết quả, tất cả các câu hỏi đều ở mức nhu cầu rất cao. Đáng chú ý, ở Q6, 100\% người trả lời đồng ý rằng sinh viên cơ khí phải có khả năng chống chịu cao để đối mặt với nhiều tình huống có thể xảy ra, đồng thời phải có trách nhiệm với sự tin cậy được giao phó. Ở Q5, 96\% người trả lời đồng ý rằng sinh viên ngành này cần có khả năng đưa ra quyết định dựa trên bằng chứng xác thực. Trong khi đó, 94\% cho rằng sinh viên cơ khí cần có khả năng nhận diện và phân tích mọi vấn đề trong mọi tình huống, và phát triển kỹ năng tư duy (lần lượt ứng với Q1 và Q2). Sinh viên phải khéo léo trong việc diễn giải kết quả đánh giá và trình bày rõ ràng lời giải thích hoặc ý tưởng. Khả năng này có thể được rèn luyện và áp dụng trong các bài đánh giá và thảo luận nhóm.

\textbf{Bảng 4.} Kết quả của bảng hỏi bảy mục.

{\small
\begin{longtable}[]{@{}lll@{}}
\toprule\noalign{}
Câu hỏi & Mức nhu cầu & Tỷ lệ phần trăm \\
\midrule\noalign{}
\endhead
\bottomrule\noalign{}
\endlastfoot
Q1 & Rất cao & 94,00 \\
Q2 & Rất cao & 94,00 \\
Q3 & Rất cao & 93,66 \\
Q4 & Rất cao & 86,33 \\
Q5 & Rất cao & 95,66 \\
Q6 & Rất cao & 100,00 \\
Q7 & Rất cao & 91,33 \\
\end{longtable}
}

Kết quả của Q3 cũng cho thấy gần 94\% kỹ sư cơ khí cần khéo léo trong việc tạo ra ý tưởng mới. Họ cần sáng tạo hơn và nắm bắt các vấn đề/công nghệ hiện tại để tạo ý tưởng thay thế nhằm hoàn thành nhiệm vụ. Kết quả 91\% ở Q7 cho thấy sinh viên ngành cơ khí phải linh hoạt đối mặt và thích nghi với mọi tình huống để hiểu văn hóa làm việc dù ở một nơi mới. Cuối cùng, ở Q4, 86\% người trả lời đồng ý rằng tư duy vượt ra ngoài lối mòn là quan trọng đối với mọi sinh viên cơ khí để bảo đảm mọi hành động được thực hiện chín chắn hơn và vượt ngoài kỳ vọng.

\subsection{3.3 Đánh giá phản hồi từ bảng hỏi}

Bảng 5 trình bày cách phân mức theo điểm trung bình được dùng trong nghiên cứu. Nhà nghiên cứu dùng các ngưỡng điểm trung bình: i. 3,68 - 5,00 (nhu cầu cao); ii. 2,34 - 3,67 (nhu cầu trung bình); iii. 1,00 - 2,33 (nhu cầu thấp). Phương pháp này giúp đơn giản hóa việc phân loại kết quả.

\textbf{Bảng 5.} Phân loại mức yêu cầu theo điểm trung bình.

Công thức: (5 - 1) / 3 = 1,33

{\small
\begin{longtable}[]{@{}lll@{}}
\toprule\noalign{}
Điểm trung bình & Mức & Diễn giải \\
\midrule\noalign{}
\endhead
\bottomrule\noalign{}
\endlastfoot
3,68 - 5,00 & Ba & Cao \\
2,34 - 3,67 & Hai & Trung bình \\
1,00 - 2,33 & Một & Thấp \\
\end{longtable}
}

(Nguồn: điều chỉnh từ thang đo Likert, 1932)

Hình 3 thể hiện biểu đồ hộp (box plot) của điểm trung bình Kỹ năng Tư duy phản biện và Giải quyết vấn đề, được tạo từ các câu hỏi mà người trả lời đã trả lời. Thông tin từ hình 3 được chuyển vào Bảng 6 và toàn bộ kết quả được thể hiện rõ.

\textbf{Hình 3.} Biểu đồ hộp điểm trung bình Kỹ năng Tư duy phản biện và Giải quyết vấn đề.

Bảng 6 trình bày điểm trung bình và thống kê tóm tắt của Kỹ năng Tư duy phản biện và Giải quyết vấn đề. Theo kết quả, mỗi người trả lời đều trả lời toàn bộ bảng hỏi với điểm tối thiểu 3,00 và tối đa 5,00. Tương tự, không người trả lời nào chọn mức 1,00 và 2,00 cho bất kỳ mục nào. Mục có điểm trung bình cao nhất là Q3 (4,57), tiếp theo là Q6 (4,56), Q1 (4,52), Q7 (4,48), Q2 (4,47), Q5 (4,43), và thấp nhất là Q4 (4,28). Kết quả đáp ứng yêu cầu và phù hợp với Bảng 4, trong đó điểm trung bình 3,68 - 5,00 thuộc mức cao. Do đó, điểm trung bình cho thấy cả bảy mục đều ở mức nhu cầu cao, và điều này cũng được Bảng 6 và các tài liệu tham khảo {[}28-31{]} củng cố.

\textbf{Bảng 6.} Thống kê tóm tắt Kỹ năng Tư duy phản biện và Giải quyết vấn đề.

{\small
\begin{longtable}[]{@{}
  >{\raggedright\arraybackslash}p{(\columnwidth - 8\tabcolsep) * \real{0.1940}}
  >{\raggedright\arraybackslash}p{(\columnwidth - 8\tabcolsep) * \real{0.2687}}
  >{\raggedright\arraybackslash}p{(\columnwidth - 8\tabcolsep) * \real{0.1642}}
  >{\raggedright\arraybackslash}p{(\columnwidth - 8\tabcolsep) * \real{0.1194}}
  >{\raggedright\arraybackslash}p{(\columnwidth - 8\tabcolsep) * \real{0.2537}}@{}}
\toprule\noalign{}
\begin{minipage}[b]{\linewidth}\raggedright
Mục câu hỏi
\end{minipage} & \begin{minipage}[b]{\linewidth}\raggedright
Số người trả lời
\end{minipage} & \begin{minipage}[b]{\linewidth}\raggedright
Tối thiểu
\end{minipage} & \begin{minipage}[b]{\linewidth}\raggedright
Tối đa
\end{minipage} & \begin{minipage}[b]{\linewidth}\raggedright
Điểm trung bình
\end{minipage} \\
\midrule\noalign{}
\endhead
\bottomrule\noalign{}
\endlastfoot
Q1 & 300 & 3,00 & 5,00 & 4,52 \\
Q2 & 300 & 3,00 & 5,00 & 4,47 \\
Q3 & 300 & 3,00 & 5,00 & 4,57 \\
Q4 & 300 & 3,00 & 5,00 & 4,28 \\
Q5 & 300 & 3,00 & 5,00 & 4,43 \\
Q6 & 300 & 4,00 & 5,00 & 4,56 \\
Q7 & 300 & 3,00 & 5,00 & 4,48 \\
\end{longtable}
}

Dựa trên Phân tích mô tả (DA), tỷ lệ phần trăm của toàn bộ kết quả được xét theo mức "cần" và "rất cần". Chỉ có hai mục cho thấy nhu cầu, dù là yêu cầu cao hay thấp của ngành công nghiệp trong việc đào tạo kỹ sư cơ khí có năng lực cạnh tranh; tỷ lệ phần trăm được thể hiện ở Bảng 7.

\textbf{Bảng 7.} Phân loại tỷ lệ phần trăm theo mức yêu cầu.

Công thức: (100 - 1) / 5 = 19,8

{\small
\begin{longtable}[]{@{}lll@{}}
\toprule\noalign{}
Tỷ lệ phần trăm & Mức & Diễn giải \\
\midrule\noalign{}
\endhead
\bottomrule\noalign{}
\endlastfoot
80,6 - 100 & Năm & Rất cao \\
60,7 - 80,5 & Bốn & Cao \\
40,8 - 60,6 & Ba & Trung bình \\
20,9 - 40,7 & Hai & Thấp \\
1,00 - 20,8 & Một & Rất thấp \\
\end{longtable}
}

(Nguồn: điều chỉnh từ thang đo Likert, 1932)

Hình 4 biểu diễn các câu hỏi ở mức yêu cầu cao. Trục hoành biểu diễn câu hỏi mà người trả lời đã trả lời, trục tung cho biết tỷ lệ phần trăm kết quả của từng mục. Cả bảy mục đều ở mức yêu cầu cao và cũng thỏa điều kiện ở Bảng 7, tức giá trị phần trăm 80,6 - 100,00 là mức rất cao. Điều này cho thấy các yếu tố của Kỹ năng Tư duy phản biện và Giải quyết vấn đề là cần thiết đối với sinh viên kỹ thuật, như đã nêu trong tài liệu tham khảo {[}18{]} và {[}33{]}.

Hình 4. Biểu đồ các câu hỏi khảo sát về Tư duy phản biện và Kỹ năng giải quyết vấn đề.

Bảng 8 trình bày tỷ lệ phần trăm tổng hợp của toàn bộ các câu hỏi về Tư duy phản biện và Kỹ năng giải quyết vấn đề. Mức 94\% cho thấy nhu cầu rất cao đối với các kỹ sư cơ khí có năng lực cạnh tranh theo nhu cầu hiện tại của ngành; các kỹ năng này cần được nắm vững hoàn toàn trước khi bước vào thực tế, theo tài liệu {[}39{]}.

Bảng 8. Tổng hợp tỷ lệ phần trăm chung.

{\small
\begin{longtable}[]{@{}lll@{}}
\toprule\noalign{}
Hạng mục & Mức nhu cầu & Giá trị \\
\midrule\noalign{}
\endhead
\bottomrule\noalign{}
\endlastfoot
Phần trăm & Rất cao & 93,57 \\
\end{longtable}
}

Theo các tài liệu {[}38,40{]}, các kỹ năng chung là kim chỉ nam giúp mỗi kỹ sư tự đánh giá mức độ kiến thức, kỹ năng và thái độ đối với nghề; hơn thế nữa, ngày nay ngành công nghiệp còn cần thêm các kỹ năng bổ sung, trong đó có tư duy phản biện và khả năng giải quyết vấn đề của kỹ sư Malaysia để cạnh tranh với các nước khác trên thế giới.

\subsection{3.4 Củng cố thành phần tư duy phản biện}

Để đào tạo kỹ sư chuyên nghiệp vững vàng, nhiều cách tiếp cận và phương pháp được áp dụng trong các học phần kỹ thuật ở bậc đại học có liên quan đến các thành phần của tư duy phản biện. Có năm thành phần cơ bản của tư duy phản biện cần thiết có thể đưa vào các môn học của sinh viên ngành Kỹ thuật cơ khí: kiến thức chuyên môn, kinh nghiệm, năng lực, thái độ và chuẩn mực nghề nghiệp. Sự kết hợp năm thành phần này tạo ra đầu ra kỹ sư có năng lực, như đề xuất ở {[}41,42{]}. Nội dung của chương trình đào tạo kỹ thuật có liên quan đến văn bản và thang chấm điểm (rubric) cần được rà soát theo phạm vi tư duy phản biện và giải quyết vấn đề. Trong thời gian học đại học, sinh viên kỹ thuật phải biết trình bày mạch lạc, sáng tạo và hợp lý với mức độ tư duy phản biện cao; đồng thời, khi giải quyết bất kỳ vấn đề nào, phải biết cân nhắc đồng thời nhiều khái niệm để đi đến giải pháp khả thi.

Theo Kế hoạch Malaysia lần thứ 11 giai đoạn 2016 - 2020: Lộ trình (10.77), mọi cơ sở giáo dục đại học phải thường xuyên rà soát các chương trình đang mở, loại bỏ chương trình trùng lặp, không còn phù hợp với ngành và không mang lại lợi ích. Nỗ lực liên tục nâng cao chất lượng và mở chương trình mới sẽ thông qua rà soát và sàng lọc nghiêm ngặt nhằm bảo đảm chương trình có giáo trình tốt hơn, giá trị gia tăng phù hợp và đáp ứng nhu cầu của ngành. Theo nghiên cứu của {[}42{]}, cách tiếp cận tư duy phản biện và giải quyết vấn đề cần được giảng dạy tại nhiều cơ sở học thuật, không chỉ ở đại học công lập mà cả đại học tư thục có đào tạo kỹ thuật.

Nghiên cứu này nhấn mạnh tổng quan về các yếu tố quan trọng của kỹ năng tư duy và giải quyết vấn đề được áp dụng trong chương trình giáo dục kỹ thuật. Theo trang JobStreet.com, thiếu kỹ năng tư duy phản biện và giải quyết vấn đề được báo cáo phổ biến là thiếu hụt kỹ năng thiết yếu ở 45\% ứng viên, một tình trạng đáng lo ngại đối với thế hệ kỹ sư tương lai.

\section{4. Kết luận}

Để đào tạo những sinh viên cơ khí có tiềm năng cao, đáp ứng khoảng trống của thị trường lao động hiện nay, cần nhấn mạnh các yếu tố của tư duy phản biện và kỹ năng giải quyết vấn đề. Các kỹ năng này nay đã trở thành mối quan tâm thực sự chứ không còn chỉ là một yêu cầu. Sinh viên ngành này được khuyên không nên chỉ dựa vào giáo trình học thuật, vì điều đó sẽ hạn chế mức độ tư duy sáng tạo. Vì vậy, giảng viên và các nhà học thuật cũng phải đóng vai trò người hỗ trợ, hướng dẫn và thực hành quá trình tự phản tư (self-reflection) vốn là yếu tố thiết yếu của tư duy phản biện, đồng thời từ bỏ lối dạy "thầy đọc trò chép" (chalk and talk). Nghiên cứu cũng chứng minh rằng đây là yếu tố kỹ năng mềm mà ngành công nghiệp cần để có được những kỹ sư chất lượng cao và tích cực.

Môi trường ngành Kỹ thuật Cơ khí luôn vận động và tiến lên theo sự phát triển nhanh chóng của công nghệ, nên không quá lời khi nói rằng sinh viên tốt nghiệp và kỹ sư cơ khí cần chú ý hơn đến việc vận dụng và khám phá tư duy phản biện trong sự nghiệp tương lai. Do đó, năng lực kỹ năng mềm là tiêu chí chính mà doanh nghiệp yêu cầu khi tuyển dụng những nhân viên thành thạo chuyên môn và những kỹ sư toàn diện.

References {[}1{]} Gagne R 1988 Instructional Science. 17(4) 387-390 {[}2{]} Halpern DF 1989 Thought and knowledge: An introduction to critical thinking (Hillsdale, NJ: Lawrence Erlbaum) {[}3{]} Ennis RH 1991 Goals for a critical thinking curriculum ed AL Costa {[}4{]} Snyder LG and Snyder MJ 2008 Delta Pi Epsilon Journal L(2) 90-99 {[}5{]} Halpern DF 1998 American Psychologist 53(4) 449-455 {[}6{]} Dantas-Whitney M 2002 System 30(4) 543-555 {[}7{]} Jeevanantham LS 2005 Why teach critical thinking? Africa Education Review 2(1) 118-129 {[}8{]} William D, Hendricson MA, Sandra MS, Andrieu C, Chadwick DG, Jacqueline E, Chmar BA, James R, Cole, Mary C, George RDH, Gerald N, Glickman, Joel FJD, Glover, Jerold S, Goldberg, Haden NK, Cyril Meyerowitz, Neumann L, Pyle M, Lisa A, Tedesco, Richard W, Valachovic MPH, Richard G, Weaver, Ronald L, Winder, Young SK, Kenneth L and Kalkwarf 2006 J. Dental Education 70(9) 925-936 {[}9{]} Kathpalia SS and Heah C 2008 RELC J. 39 300-317 {[}10{]} Reinstein-Staton R 2008 The rights and wrongs of strategic planning', security: for buyers of products, systems \& services. 45(7) 34-36 {[}11{]} Nugent PM and Vitale BA 2008 Fundamental success: a course review applying critical thinking to test taking, USA: F. A. Davis Company {[}12{]} Lau J 2012 A Mini Guide to Critical Thinking {[}13{]} Stevens R, Johri A and O'Connor K 2014 Professional engineering work ed Johri A and Olds BM (Chapter 7: Cambridge handbook of engineering education research, Cambridge University Press, New York) {[}14{]} Mohamad S, Rasul R, Amnah AR and Mohd YH 2014 Kemahiran kebolehdapatan kerja suatu keperluan pekerjaan (Penerbit Universiti Kebangsaan Malaysia, Bangi) {[}15{]} Kamarudin MT, Ruhizan MY and Ramlee M 2014 Pemangkin kebolehpasaran graduan, status pekerjaan graduan, 30 (Penerbit Universiti Kebangsaan Malaysia, Bangi) {[}16{]} Mohd Huzairi J, Yuzainee MY, Mohd Zaidi O and Azami Z 2014 Kemahiran kebolehgajian graduan kejuruteraan (Penerbit Universiti Kebangsaan Malaysia, Bangi) {[}17{]} Aldridge MD 1994 J. Eng. Education 83(3) 231-236 {[}18{]} Eraut M 1994 Developing professional knowledge and competence (Falmer Press, London) {[}19{]} Walker SE 2003 J. Athletic Training 38(3) 263-267 {[}20{]} Nickerson RS 1994 The teaching of thinking and problem-solving ed Sternberg RJ Thinking and problem-solving (2nd ed. San Diego: Academic Press) {[}21{]} Leach BT and Good DW 2011 Int. J. Humanities and Social Science 1(21) 100-106 {[}22{]} Aliakbari M and Sadeghdaghighi A 2011 Investigation of the relationship between gender, field of study, and critical thinking skill: the case of Iranian students. Paper Presented at The 16th Conference of Pan-Pacific Association of Applied Linguistics, The Chinese University of Hong Kong 301-310 {[}23{]} Mahdyeh N and Arefi MA 2014 Indian J. Fundamental and Applied Life Sciences 4(1) 153-162 {[}24{]} Funke J 2013 J. Problem Solving 6(1) 2-19 {[}25{]} Hasnah I and Jamaludin B 2017 Malay Language Education J. 7(1) 56-65 {[}26{]} Vanderbilt University School of Engineering {[}27{]} Rosdiadee N 2013 Technical communication skills among recent electrical and electronics engineering graduates in job industries 164 (Universiti Kebangsaan Malaysia, Bangi, Selangor, Malaysia) {[}28{]} Pusat Perkembangan Kurikulum 2002 Kemahiran berfikir dalam pengajaran dan pembelajaran (Pusat Perkembangan Kurikulum, Kementerian Pendidikan Malaysia) {[}29{]} Bahagian Jaminan Kualiti 2004 Panduan kriteria dan standard bagi program bidang pendidikan (Jabatan Pendidikan Tinggi, Kementerian Pendidikan Malaysia) {[}30{]} Roselainy AR, Yudariah MY, Hamidreza K and Sabariah B 2012 Developing mathematical communication skills of engineering students (Procedia-Social and Behavioral Science) {[}31{]} Collins JP 2008 Modern approaches to teaching and learning anatomy. Student BMJ 337:a1753 {[}32{[} Colley BM, Bilics AR and Lerch CM 2012 The Canadian J. for the Scholarship of Teaching and Learning 3(1) 1-19 {[}33{]} Jake M and Laguador 2013 Developing students' attitude leading towards a life-changing Career (Educational research international) {[}34{]} Infor Resources System Sdn. Bhd. 2001 Creativity and creative problem-solving (IRSSB Selangor) {[}35{]} Ili Atiqah Md D and Ruslin A 2016 J. Psikologi Malaysia 30(1) 16,

http://www.wolcottlynch.com {[}36{]} Campos HM, Rubio AM, Atondo GH, Palma YM and Chorres 2015 Relationship between creativity, personality, and entrepreneurship: An exploratory study (Mexico) {[}37{]} Mishra DS 2016 Int. Research J. Eng. and Tech. 3(2) {[}38{]} Inayatullah K, Noor Abidah MO, Yusuf B and Zafar Iqbal SM 2012 Int. J. Applied Linguistics \& English Literature 1(5) 176-183 {[}39{]} Anderson C and John FG 2013 Skills requirement for tomorrow's best jobs helping educators provide students with skills and tools they need (IDC Analyse the Future) {[}40{]} Desmond A and Martin J 2016 Int. J. Higher Education 5(2) 23-59 {[}41{]} Hairuzila I, Hazadiah MD and Normah A 2010 2nd Int. Congress on Engineering Education, Kuala Lumpur, Malaysia 258-263 Lời cảm ơn Nghiên cứu được tài trợ toàn bộ bởi Khoa Thiết kế và Công nghệ Đổi mới, Universiti Sultan Zainal Abidin (UniSZA), Trung tâm Nghiên cứu Giáo dục Kỹ thuật, Khoa Kỹ thuật và Môi trường Xây dựng, Đại học Quốc gia Malaysia (UKM), và Bộ Giáo dục Đại học (MoHE).
